\begin{lemma}[The middle Rutgers exit lies beyond the terminal ladder]
Fix an original positive-rung extended Tashkinov tree \(\mathcal E\).
Assume \(\Delta(G)+1\leq k\).  Assume current assertion (iv) uniformly
for every maximum ETT having the same rung number as \(\mathcal E\): at a
terminal PE, under every proper terminal-ladder-stable coloring leaving
exactly \(e_0\) uncolored, the terminal \(\gamma_n,\delta_n\)-component
through the certified support meets the terminal ladder only at that
support.  Assume also assertion (ii) for every maximum positive-rung ETT
whose rung number is strictly smaller than that of \(\mathcal E\), with
arbitrary certified input.

Let \(\mathcal P=(\sigma,U,\ldots)\) be a
\noderef{RutgersClosedPresentation} relative to \(\mathcal E\), let
\(\mathcal C\) be a \noderef{RutgersCounterexample} on \(\mathcal P\),
and suppose that \(V(U)\) is elementary under \(\sigma\).  Write
\(\mathcal R=\mathcal P.\mathit{corresponding}.\mathit{replacementETT}\),
write \(n\) for its positive rung number, and write \(T_n\) for its
actual terminal ladder tree.  Then the middle ordered tree end of
\(\mathcal C\) satisfies
\[
                         b_2\notin V(T_n).
\]
No minimum property of \(\mathcal C\) is assumed.
\end{lemma}
\begin{proof}
Put \(D_n=\noderef{TashkinovStoredColors}(\mathcal R,n)\), with \(n\)
denoting the terminal data index as well as the rung number when no
confusion can result.  The final-generating-coloring and dependent-series
fields of \noderef{CorrespondingExtendedTashkinovTree} identify
\(\sigma\) with the terminal coloring in the actual certified series of
\(\mathcal R\).  Eliminating its positive terminal certificate in
\noderef{TashkinovSeries} and \noderef{TashkinovSeriesStep} shows that
this coloring is proper in the sense of
\noderef{IsProperEdgeColoring} and leaves exactly \(e_0\) uncolored in
the sense of \noderef{IsOneEdgeUncolored}: these are explicit fields in
SE and PE, and in RE they follow from the preceding proper coloring and
the current-coloring equality.  We use these facts below without changing
from the replacement presentation to the original ETT.  All references
below to the fields of \(\mathcal C\), including its colors, walks,
endpoints, support-disjointness, and retained boundary identities, are to
the definition \noderef{RutgersCounterexample}.

Suppose, for a contradiction, that \(b_2\in V(T_n)\).  We first record
the exact prefix relation
\[
                              T_n\subseteq U.                 \tag{1}
\]
In the PE alternative of the terminal route stored in
\noderef{RutgersClosedPresentation},
\(\noderef{TashkinovClosure}(\sigma,T_n,U)\) contains a
\noderef{TashkinovAugmentationSequence} from \(T_n\) to \(U\), so
\noderef{TashkinovAugmentationSequenceInitialSegment} gives (1).  In the
SE/RE alternative, the route contains an exact seed \(A\), an exact edge
identity \(f\), and
\(\noderef{TreeSequenceAdjoinEdge}(T_n,f,A)\), followed by a closure
from \(A\) to \(U\).  The adjoin-edge definition makes \(T_n\) an exact
\noderef{TreeSequenceInitialSegment} of \(A\), and the same augmentation
lemma makes \(A\) an exact initial segment of \(U\).  Composing their
length inequalities and their entrywise vertex and edge equalities gives
(1).  The vertex sets here and below are those of
\noderef{TreeSequenceVertexSet}.  In either alternative, the closure
field also says that \(V(U)\) is closed under \(\sigma\), in the precise
sense of \noderef{ColoredVertexSetClosed}.

Expand \(b_1\prec_U b_2\) from the field
\(\mathcal C.\mathit{insideOrder12}\) using
\noderef{TreeVertexOrder}.  It supplies positions \(i<j\) of \(b_1,b_2\)
in \(U\).  Membership of \(b_2\) in \(T_n\), (1), and injectivity of
the vertex array in the underlying tree-sequence identify \(j\) with the
image of its position in \(T_n\).  Hence \(j\leq |T_n|\), and then
\(i<j\) says that position \(i\) also lies in the exact prefix.  The
entrywise equality in \noderef{TreeSequenceInitialSegment} yields
\[
                         b_1,b_2\in V(T_n).                   \tag{2}
\]
The vertex-set form of (1), equivalently
\noderef{TreeSequenceInitialSegmentVertexSetSubset}, restricts the
assumed \noderef{ColoredVertexSetElementary} relation on \(V(U)\) to
\(V(T_n)\).

The first clause of \noderef{IsOneEdgeUncolored} gives
\(\sigma(e_0)=\mathsf{none}\).
Apply \noderef{TreeMissingColorsCardLowerBound} to the elementary set
\(V(T_n)\), this equality, and \(\Delta(G)+1\leq k\).  It gives
\[
             |V(T_n)|+2\leq |\overline\sigma(T_n)|.          \tag{3}
\]
In particular \(\overline\sigma(T_n)\) is nonempty.  In an SE or RE
terminal step choose any
\(c\in\overline\sigma(T_n)\).  In a PE step first expose from the actual
positive constructor its certified support \(v_n\), colors
\(\gamma_n,\delta_n\), and exact connecting edge \(f_n\).  A
tree-sequence in the sense of \noderef{TreeSequence} contains the two
distinct ends of its first edge; their distinctness is the loopless law
of \noderef{LooplessMultigraph}.  Thus \(|V(T_n)|\geq2\), and (3) gives
at least four missing colors.  Choose
\[
 c\in\overline\sigma(T_n)-\{\gamma_n,\delta_n\}.             \tag{4}
\]
This slightly more selective choice removes any ambiguity when the PE
switch is later undone.

By (1), the chosen color \(c\) is also missing on \(U\), while
\(\alpha=\mathcal C.\alpha\) is missing on \(U\) by the counterexample
field.  Here and below tree-missingness is the union in
\noderef{TreeMissingColors} of the vertex missing-color sets from
\noderef{MissingColors}.  Define \(\mu_1=\sigma\) if \(c=\alpha\).
Otherwise put
\[
 \mu_1=\noderef{OutsideTreeColorSwap}(U,\sigma,\alpha,c).
\]
In the second case the two colors are distinct, \(U\) is closed, and
both are missing on \(U\), so
\noderef{OutsideTreeColorSwapStable}, with control set \(D_n\), proves
that \(\mu_1\) is proper, leaves exactly \(e_0\) uncolored, is
\((U,D_n,\sigma)\)-stable, and fixes every exact edge identity incident
with \(U\).  In the first case the same conclusions are the reflexive
case of \noderef{StableColoring}, also supplied by the reflexivity clause
of \noderef{sc2}.  Restricting the vertex equality and incident-edge
clauses along (1) gives
\[
            \mu_1\text{ is }(T_n,D_n,\sigma)\text{-stable}.  \tag{5}
\]
The vertexwise missing-set equality in \noderef{StableColoring}, followed
by the union definition of \noderef{TreeMissingColors}, also gives
\[
                         c\in\overline{\mu}_1(T_n).            \tag{5a}
\]

We next verify from the stored walks, rather than assume, what happens to
the first two exit legs.  The field
\(\mathcal C.\mathit{betaPresent}\) says that \(\beta\) is missing at no
vertex of \(U\); since \(c\) is missing at a vertex of \(T_n\subseteq U\),
we have \(c\ne\beta\).  Fix \(i\in\{1,2\}\).  If \(c=\alpha\), the
original \(U\)-exit certificate restricts immediately to \(T_n\) by
(1)--(2).  Assume henceforth in this paragraph that \(c\ne\alpha\).

First we rebuild the incidence-walk clause.  Let \(e\) be an edge-node
of \(Q_i\).  If \(e\) is not incident with \(U\) in the sense of
\noderef{TreeIncidentEdge}, then
\noderef{OutsideTreeColorSwap} changes its old color \(\alpha\) to
\(c\) and leaves its old color \(\beta\) unchanged.  Suppose instead
that \(e\) is incident with \(U\).  Since both endpoints of \(Q_i\) are
vertex-nodes, the edge-node \(e\) is internal on this simple path.  Its
predecessor and successor are distinct vertex-nodes, by path simplicity,
and adjacency in \noderef{TwoColorIncidenceGraph} says that both are
ends of \(e\).  The endpoint-pair description and loopless law in
\noderef{LooplessMultigraph} say that these are precisely the two
distinct ends of \(e\).  At least one is in \(V(U)\) by
\noderef{TreeIncidentEdge}; the unique-contact clause of the original
\noderef{TreeExitPath} forces that end to be \(b_i\) and forces the other
end outside \(V(U)\).  Thus \(e\) is an exact boundary identity in
\noderef{BoundaryEdges}\((G,V(U))\).

Both \(\alpha\) and \(c\) are missing on \(U\).  Choosing their
respective missing vertices by \noderef{TreeMissingColors} and applying
the closedness predicate \noderef{ColoredVertexSetClosed} to this exact
boundary identity shows that \(\sigma(e)\) is neither \(\alpha\) nor
\(c\).  But the old walk adjacency says that \(\sigma(e)\) is
\(\alpha\) or \(\beta\), so it is \(\beta\).  The outside swap fixes
this tree-incident identity.  We have now checked every edge-node, and
the unchanged node sequence \(Q_i\) is therefore a walk in
\noderef{TwoColorIncidenceGraph}\((G,\mu_1;c,\beta)\).

Its path and length clauses are unchanged.  Its outside endpoint \(a_i\)
is outside \(U\), hence outside \(T_n\); (2) puts \(b_i\) in \(T_n\).
If a vertex-node of \(Q_i\) belongs to \(V(T_n)\), then (1) puts it in
\(V(U)\), so the original unique-contact clause identifies it with
\(b_i\).  It remains only to verify the missing-color alternative at
\(a_i\).  If \(\beta\) was missing there under \(\sigma\), it remains
missing: \(\beta\ne\alpha,c\), and the outside swap never creates or
removes a \(\beta\)-colored edge.  Suppose instead that \(\alpha\) was
missing there.  If some edge \(g\) incident with \(a_i\) had color
\(c\) under \(\mu_1\), split according to
\noderef{TreeIncidentEdge}\((U,g)\).  If \(g\) is not tree-incident,
the definition of the swap and \(c\ne\alpha\) give
\(\sigma(g)=\alpha\), contrary to \(\alpha\)-missingness.  If \(g\)
is tree-incident, it is fixed and has old color \(c\).  Since
\(a_i\notin V(U)\), its incidence with \(a_i\) and with a vertex of
\(U\), together with the loopless endpoint description, put \(g\) in
\noderef{BoundaryEdges}\((G,V(U))\).  This contradicts closedness of
\(U\), because \(c\) is missing on \(U\).  Thus no such \(g\) exists.
By the edge-witness definition \noderef{PresentColors} and its complement
\noderef{MissingColors}, \(c\) is missing at \(a_i\) under \(\mu_1\).

These verifications give every clause of \noderef{TreeExitPath} for
both \(i=1,2\).  Their supports are still disjoint because the stored
node sequences have not changed.  Consequently
\[
 Q_1,Q_2\text{ are support-disjoint }
       (T_n,\mu_1,\{c,\beta\})\text{-exit paths}.             \tag{6}
\]
The exact identities \(f_1,f_2\) remain on their respective walks and
remain incident with \(b_1,b_2\).  Expanding their old membership in
\noderef{BoundaryColorEdges}\((G,\sigma,V(U),\beta)\) says that each is
an exact \noderef{BoundaryEdges}-identity of \(U\) colored \(\beta\).
Since its inside endpoint \(b_i\) lies in \(T_n\) by (2), while its
other endpoint lies outside \(U\) and hence outside \(T_n\), the same
identity lies in the \(\beta\)-colored boundary of \(T_n\) under
\(\sigma\).  It is incident with \(U\), so the outside swap fixes it and
it remains in that boundary under \(\mu_1\).  Finally \(f_1\ne f_2\):
equality would put the same edge-node in both walk supports, contrary to
the recorded support-disjointness.

We isolate a consequence of (6) which will be used after either terminal
transport:
\[
 \begin{gathered}
 \rho\text{ is proper},\quad c\ne\beta,\quad
 c\in\overline\rho(T_n),\quad a_1\ne a_2,\\
 \noderef{TreeExitPath}(T_n,\rho,c,\beta,a_i,b_i,Q_i)
 \text{ holds exactly for }i=1,2\\
 \Longrightarrow\quad
 \neg\noderef{TreeInterchangeable}(G,T_n,\rho,c,\beta).
                                                               \tag{7}
 \end{gathered}
\]
We prove (7), including the possibility of isolated components.  First
observe that whenever a vertex \(w\) misses at least one of \(c,\beta\)
under \(\rho\), its vertex-node has incidence degree at most one.  Indeed,
by \noderef{TwoColorIncidenceGraph} every neighbor is an exact edge
identity incident with \(w\) and colored \(c\) or \(\beta\).  If \(c\)
is missing, the complement relation between \noderef{MissingColors} and
\noderef{PresentColors} rules out the first color, and properness in
\noderef{IsProperEdgeColoring} allows at most one incident edge of color
\(\beta\).  If \(\beta\) is missing, the symmetric argument allows at
most one edge of color \(c\).  This proves the claimed bound even when
the neighbor set is empty.

Choose \(x\in V(T_n)\) missing \(c\), using the union definition
\noderef{TreeMissingColors}.  The preceding degree bound applies to
\(x\), and the missing-color clause in each exact exit certificate in
(7) applies it to \(a_1,a_2\).  Because \(c\ne\beta\), the root itself
belongs to its full connected component and has a subsingleton neighbor
set.  These are exactly the clauses of
\noderef{TwoColorPathComponent} for the components rooted at
\(x,a_1,a_2\); no positive-degree assumption has been made.  The first
component meets \(T_n\) at \(x\).  For \(i=1,2\), the walk \(Q_i\)
connects \(a_i\) to \(b_i\in V(T_n)\), so equality of connected
components in \noderef{TwoColorComponent} shows that the component rooted
at \(a_i\) meets \(T_n\) at \(b_i\).

If \(c,\beta\) were \(T_n\)-interchangeable, the defining implication in
\noderef{TreeInterchangeable}, first for the roots \(x,a_1\) and then for
\(a_1,a_2\), would identify all three full components.  Apply
\noderef{TwoColorComponentSpanningPath} to this common component; its
hypothesis is precisely the assumed properness of \(\rho\).  The result
is a finite simple path whose support is the whole component.  Each of the
three vertex-nodes lies on this path and has ambient incidence degree at
most one.  Such a node must be a path endpoint: if it occurred at an
internal index, its predecessor and successor would be distinct by path
simplicity and would give two distinct incidence-graph neighbors.  But
\(x,a_1,a_2\) are pairwise distinct, since \(x\in V(T_n)\), each
\(a_i\notin V(T_n)\), and \(a_1\ne a_2\).  A path has only its first
and last node as endpoints, so three distinct nodes cannot all be
endpoints.  This contradiction proves (7).  In particular, the argument
uses the at-most-one bound rather than a degree-one assertion and thus
includes the isolated-component convention built into
\noderef{TwoColorPathComponent}.

We now transport to the preceding rung, splitting the actual terminal
constructor.

If it is SE, its \noderef{MaximumDefectiveChoice} field gives a selected
coloring which is \((T_n,D_{n-1},\varphi_{n-1})\)-stable, and the current
coloring equation identifies that selected coloring with \(\sigma\).
The first inclusion of \noderef{HkuMissingStoredMonotone} says
\[
 \overline{\varphi}_{n-1}(T_n)\cup D_{n-1}
 \subseteq \overline\sigma(T_n)\cup D_n.                    \tag{8}
\]
Using (8) in the incident-edge clause of (5), while restricting its
vertex clause, makes \(\mu_1\)
\((T_n,D_{n-1},\sigma)\)-stable.  Transitivity in
\noderef{sc2} with the selected-stable certificate gives
\[
 \mu_1\text{ is }(T_n,D_{n-1},\varphi_{n-1})\text{-stable}.  \tag{9}
\]
In RE, the constructor equation is \(\sigma=\varphi_{n-1}\); the same
use of (8) gives (9) directly.  Properness and exact-one-uncoloredness of
every coloring used in this application are fields of the actual
constructor or were proved above.

Suppose instead that the terminal kind is PE.  Use its actual constructor
witnesses \(v_n,\gamma_n,\delta_n,f_n\).  The
\noderef{ParallelExtensionColor} field says that \(\gamma_n\) is missing
at \(v_n\) before the terminal switch.  The exact selected edge from
\noderef{MaximumDefectiveChoice} is incident with \(v_n\) and has color
\(\delta_n\); the intermediate stable relation preserves that value.
Thus \(\gamma_n\ne\delta_n\).  Expanding
\noderef{TwoColorComponentSwap}, properness as defined in
\noderef{IsProperEdgeColoring} makes that edge the unique
\(\delta_n\)-edge at \(v_n\), so the terminal switch makes
\(\delta_n\) missing at \(v_n\) under \(\sigma\).  Since
\(v_n\in V(T_n)\subseteq V(U)\), the global nonmissingness of \(\beta\)
on \(U\) gives
\[
                             \beta\ne\delta_n.                \tag{10}
\]

The unconditional first conclusion of
\noderef{HkuParallelBoundaryStructure}, applied to these exact terminal
option slots and connecting identity, is
\[
           \partial_{\sigma,\gamma_n}(T_n)=\{f_n\}.          \tag{11}
\]
If \(\beta=\gamma_n\), then the two distinct exact identities
\(f_1,f_2\) exhibited after (6) both lie in the singleton (11), a
contradiction.  Hence \(\beta\ne\gamma_n\).  Together with the PE choice
(4), we have
\[
                    \{c,\beta\}\cap
                    \{\gamma_n,\delta_n\}=\varnothing.       \tag{12}
\]

Apply the assumed current assertion (iv) to the replacement ETT
\(\mathcal R\), using its rung agreement with \(\mathcal E\), the
maximum property \(\mathcal P.\mathit{maximum}\), the actual terminal
kind and option-slot equations, and (5).  It says that the
\(\gamma_n,\delta_n\)-component of \(v_n\) under \(\mu_1\) meets
\(T_n\) exactly at \(v_n\).  These are precisely the hypotheses of
\noderef{ParallelExtensionUndoStable}.  Put
\[
 \mu_2=\noderef{TwoColorComponentSwap}
          (G,\mu_1,\gamma_n,\delta_n,\operatorname{inl}v_n).
\]
The cited undo theorem gives
\[
 \mu_2\text{ is }(T_n,D_{n-1},\varphi_{n-1})\text{-stable}.  \tag{13}
\]
Because \(\gamma_n\ne\delta_n\),
\noderef{TwoColorComponentSwapProper} gives properness of \(\mu_2\).
The definition of the switch fixes \(\mathsf{none}\) and only transposes
two colored values, so the exact-one-uncolored condition for \(\mu_1\)
gives the same condition for \(\mu_2\), by the two clauses of
\noderef{IsOneEdgeUncolored}.  By (12) and the edgewise definition of
\noderef{TwoColorComponentSwap}, an exact edge has color \(c\) under
\(\mu_2\) exactly when it has color \(c\) under \(\mu_1\), and the same
equivalence holds for \(\beta\).  The adjacency definition in
\noderef{TwoColorIncidenceGraph} therefore gives identical
\(c,\beta\)-incidence graphs.  The edge-witness and complement definitions
\noderef{PresentColors} and \noderef{MissingColors} give identical
missingness for these two colors at every vertex.  Hence each exact
certificate in (6) transports field by field: the same node sequence is
a walk in the identical incidence graph; its path property, length,
support, endpoints, and contacts with \(T_n\) are unchanged; and its
outside-end missing-color alternative is preserved.  Thus \(Q_1,Q_2\)
are exact \(T_n\)-exit paths under \(\mu_2\), and the union definition
\noderef{TreeMissingColors} transports (5a) to
\(c\in\overline{\mu}_2(T_n)\).  Applying (7) gives noninterchangeability
under \(\mu_2\).

In the SE/RE cases, apply (7) directly to the exact certificates (6),
(5a), properness of \(\mu_1\), \(c\ne\beta\), and the retained
counterexample field \(a_1\ne a_2\).  In the PE case use the just
transported certificates and missingness together with properness of
\(\mu_2\).  Therefore, if \(\rho\) denotes \(\mu_1\) in the SE/RE cases
and \(\mu_2\) in the PE case, then (9) or (13), respectively, holds and
\[
             \neg\noderef{TreeInterchangeable}
                    (G,T_n,\rho,c,\beta).                    \tag{14}
\]

It remains to invoke assertion (ii) at the correct rung.  First suppose
\(1<n\).  Apply \noderef{CertifiedETTPredecessor} to \(\mathcal R\) and
\(\mathcal P.\mathit{maximum}\).  It supplies an actual positive-rung
ETT \(\mathcal R^-\) with rung number \(n-1\), maximum property,
displayed terminal tree exactly \(T_n\), generating coloring exactly
\(\varphi_{n-1}\), and terminal computed stored set exactly
\(D_{n-1}\).  Its rung is strictly smaller than the rung of
\(\mathcal E\), because
\(\mathcal P.\mathit{corresponding}.\mathit{rungAgreement}\) identifies
\(n\) with the original rung number.  Hence the assumed lower assertion
(ii), applied to \(\mathcal R^-\), the transported coloring \(\rho\),
(9) or (13), and \(c\in\overline\rho(T_n)\), says that \(c,\beta\) are
\(T_n\)-interchangeable.  This contradicts (14).

Finally suppose \(n=1\).  There is no zero-rung
\noderef{ExtendedTashkinovTree}, and none is introduced.  Under the
index equality \(n=1\), the actual tree \(T_n\), preceding coloring
\(\varphi_{n-1}\), and preceding stored set \(D_{n-1}\) are exactly the
tree, coloring, and stored set at index zero of the replacement certified
series.  The interchangeability half of
\noderef{TechnicalTheoremBase}, applied directly to that series, to
\(\rho\), to (9) or (13), and to the missing color \(c\), therefore also
says that \(c,\beta\) are \(T_n\)-interchangeable.  Again (14) is a
contradiction.  Both possible positive rung numbers have led to a
contradiction, proving \(b_2\notin V(T_n)\).
\end{proof}
